\documentclass[../main.tex]{subfiles}
\graphicspath{{\subfix{../images/}}}
\begin{document}
		\begin{itemize}
		\item Proposto agli inizi degli anni ‘70 da Codd
		\item Finalizzato alla realizzazione dell’indipendenza dei dati
		\item Unisce concetti derivati dalla teoria degli insiemi (concetto di relazione) con una rappresentazione dei dati di tipo tabellare
		\item Attualmente è largamente il modello più utilizzato
		\item Teorizzato per separare il più possibile il livello logico dal livello fisico della descrizione dei dati: la base di dati non contiene alcun riferimento all’effettiva allocazione dei dati in memoria.
		\item Basato su un rigoroso modello matematico, permette un elevato grado di astrazione.
		\item Rappresentazione tabellare: semplice ed intuitiva.\\Le relazioni ed i risultati delle operazioni sulle tabelle sono facilmente rappresentabili ed interpretabili dagli utenti.
	\end{itemize}
	\subsection{Tre accezioni importanti}
	\begin{enumerate}
		\item Relazione matematica: come nella teoria degli insiemi
		\item Relazione secondo il modello relazionale dei dati
		\item Relazione (dall’inglese \textbf{relationship}) di tipo \textbf{logico} fra entità (concetti descrivibili in modo indipendente dagli altri) nel modello Entità-Relazione (entity-relationship); traducibile anche con \textbf{associazione} o \textbf{correlazione} (scelta preferibile per evitare ambiguità).
	\end{enumerate}
	\subsection{Relazione - Prodotto Cartesiano}
	Dati due insiemi $D_1$ e $D_2$ si definisce \textbf{Prodotto Cartesiano} di $D_1$ e $D_2$ \ $(D_1 \times D_2)$ l’insieme di tutte le possibili coppie ordinate $(v_1, v_2)$ tali che $v_1$ sia un elemento di $D_1$ e $v_2$ sia un elemento di $D_2$.
	\subsubsection{Esempio}
	Dati gli insiemi: $A = \{cubo, cono\}$ e $B = \{rosso, verde, blu\}$\\
	Il prodotto cartesiano tra $A$ e $B$ è $\{(cubo, rosso), (cono, rosso), (cubo, verde), (cono, verde), (cubo, blu), (cono, blu)\}$
	\subsection{Relazioni}
	Una relazione matematica su due insiemi $D_1$ e $D_2$ è un sottoinsieme del prodotto cartesiano $D_1 \times D_2$.\\ \\
	NOTA: a livello formale gli insiemi possono essere infiniti, a livello pratico non possiamo però considerare relazioni infinite.
	\subsubsection{Esempio}
	Dati gli insiemi visti, una possibile relazione è $\{(\text{cubo}, \text{rosso}), (\text{cono}, \text{rosso}), (\text{cubo}, \text{blu})\}$\\
	o, in forma tabellare:
	\begin{tabular}{ |c|c| }
		\hline
		Cubo & Rosso\\
		Cono & Rosso\\
		Cubo & Blu\\
		\hline
	\end{tabular}
	\subsection{Tupla}
	Le definizioni viste per due insiemi possono essere generalizzate a n insiemi.\\
	Ogni riga della tabella sarà allora una n-pla ordinata di elementi (detta tupla).\\ \\
	n è detto \textbf{grado} del prodotto cartesiano e quindi della relazione.\\
	Il numero di elementi (istanze) della relazione è detto \textbf{cardinalità} della relazione.\\ \\
	Un insieme può apparire più volte in una relazione.
	\subsubsection{Perché una tupla si chiama così?}
	Perché è una n\_upla (ennupla) che viene rappresentata di solito con la lettera $t$ anziché con la lettera $n$ 
	\subsubsection{Esempio}
	La relazione $\text{Partite\_di\_Calcio} (Casa, Ospiti, RetiCasa, RetiOspiti)$ è un sottoinsieme del prodotto cartesiano \textbf{Stringa x Stringa x Intero x Intero}.
	\subsection{Cosa significa relazione}
	Una relazione è concetto mutuato dalla definizione di relazione matematica della teoria degli insiemi, definito come sottoinsieme del prodotto cartesiano fra $n$ insiemi.\\
	Nel modello relazionale corrisponde ad una struttura dati tabellare.
	\subsubsection{Proprietà fondamentali}
	\begin{itemize}
		\item Ogni tupla è internamente ordinata: l’i-esimo valore proviene dall’i-esimo dominio (struttura posizionale)
		\item Non esiste ordinamento intrinseco fra le tuple, per la natura insiemistica della relazione
		\item Non sono ammesse tuple uguali (ogni elemento di un insieme è unico)
	\end{itemize}
	\subsubsection{Conseguenze}
	\begin{itemize}
		\item Lo scambio fra righe di una tabella non modifica la relazione
		\item Lo scambio fra colonne di una tabella può portare alla sua inconsistenza con lo schema
	\end{itemize}
	\subsection{Basi di Dati e Relazioni}
	Uno \textbf{schema di relazione} $R(X)$ è costituito da un simbolo (nome della relazione) $R$ e da una serie di attributi $X = {A_1, A_2, \ldots, A_n}$\\ \\
	Quindi una tupla $t$ è una istanza di una relazione $r$ (un elemento del corrispondente insieme) la cui struttura è definita dal corrispondente schema $R(X)$.\\
	Se la relazione è rappresentata come tabella, una istanza della relazione è una riga della tabella:
	\begin{center}
		\begin{tabular}{ |c|c|c| }
			\hline
			\textbf{01} & \textbf{Analisi} & \textbf{Rossi}\\
			\hline
		\end{tabular}
	\end{center}
	Uno \textbf{schema di base di dati} è un insieme di schemi di relazione con nomi diversi
	\[ R = { R_1(X_1), R_2(X_2), \ldots, R_n(X_n) } \]
	Una \textbf{base di dati} definita su uno schema
	\[ R = { R_1(X_1), R_2(X_2), \ldots, R_n(X_n) } \]
	è un insieme di relazioni $r = { r_1, r_2, \ldots, r_n }$ dove ogni $r_i$ è una relazione sullo schema $R_i(X_i)$.\\ \\
	A livello di rappresentazione, una base di dati è un \textbf{insieme di tabelle}.
	\subsubsection{Esempio}
	$\text{Corsi}(Codice, Titolo, Docente)$\\
	Una relazione su uno schema $R(X)$ è un insieme $r$ di tuple definite su $X$: è rappresentata come una tabella:
	\begin{center}
		\begin{tabular}{ |c|c|c| }
			\hline
			\multicolumn{3}{|c|}{\textbf{Corsi}}\\
			\hline
			\textbf{Codice} & \textbf{Titolo} & \textbf{Utente}\\
			\hline
			01 & Analisi & Rossi\\
			02 & Chimica & Bruni\\
			04 & Chimica & Verdi\\
			\hline
		\end{tabular}
	\end{center}
	\subsubsection{Definizione relazione}
	Una relazione è un \textbf{insieme di record omogenei}, cioè definiti sugli stessi campi.\\
	Come ogni campo di un record è associato ad un \textbf{nome (etichetta)}, così si associa ad ogni colonna della relazione un \textbf{attributo}.
	\subsubsection{Esempio di relazione con attributi}
	\begin{center}
		\begin{tabular}{ |c|c|c|c| }
			\hline
			\multicolumn{4}{|c|}{\textbf{Corsi}}\\
			\hline
			\textbf{Casa} & \textbf{Ospiti} & \textbf{RetiCasa} & \textbf{RetiOspiti}\\
			\hline
			Parma & Inter & 3 & 2 \\
			Palermo &Lazio & 2 & 0\\
			Milan & Juventus & 1 & 1\\
			\hline
		\end{tabular}
	\end{center}
	\subsubsection{Notazione}
	Se $t$ è una tupla definita sullo schema $R(X)$ (insieme ordinato di domini) di una relazione e $A$ è uno dei domini di $R(X)$.\\
	$t[A]$ (o $t.A$) è il valore di $t$ relativo al dominio $A$.
	\subsubsection{Esempio}
	$[ \text{relazione }Partite(Casa,Ospiti,RetiCasa,RetiOspiti) ]$\\
	Se $t$ è la prima tupla della relazione (vedi pagine precedenti), allora $t.Casa = \text{Parma}$
	\begin{figure}[h]
		\centering
		\adjustbox{cfbox=white 5pt 0cm}{\includegraphics[scale=0.3]{example-image-golden}}
		\caption{Livello descrittivo e Rappresentazione}
		\label{fig:im-5}
	\end{figure}
	\subsection{Dominio}
	Ogni attributo ha un suo \textbf{dominio} su cui è definito.\\
	Quindi una tupla è un insieme di valori, uno per attributo, ordinati secondo lo schema della relazione e definiti ciascuno su un proprio dominio.\\
	Una relazione è una serie di tuple definite sullo schema della relazione (insieme ordinato dei domini dei singoli attributi).\\
	\begin{tabular}{ |l|c|}
		\hline
		Studenti($Matricola$, $Cognome$, $Nome$, $DataNascita$)\\
		Corsi($Codice$, $Titolo$, $Docente$)\\
		Esami($Studente$, $Voto$, $Corso$)\\
		\hline
	\end{tabular}
	\begin{itemize}
		\item \textbf{Studenti} contiene dati su un insieme di studenti.
		\item \textbf{Corsi} contiene dati su un insieme di corsi.
		\item \textbf{Esami} contiene dati su un insieme di esami e fa riferimento alle altre due attraverso i numeri di matricola e il codice del corso.
	\end{itemize}
	Quindi Matricola e Studente, come Corso e Codice, sono definiti sullo stesso dominio e possono (in questo caso devono, per generare il riferimento) assumere gli stessi valori.
	\begin{figure}[h]
		\centering
		\adjustbox{cfbox=white 5pt 0cm}{\includegraphics[scale=0.3]{example-image-golden}}
		\caption{Schema logico base di dati}
		\label{fig:im-6}
	\end{figure}
	Lo schema logico della base di dati è la descrizione, a livello logico, di uno schema concettuale in cui, se si usa la rappresentazione col modello Entità/Relazione:
	\begin{itemize}
		\item Studenti e Corsi sono due entità, perché sono concetti che esistono indipendentemente dagli altri e possono essere rappresentati e descritti in modo autonomo.
		\item Esami è una relazione, perché è un concetto che esiste solo come collegamento (correlazione, associazione) fra due concetti rappresentabili come entità.
	\end{itemize}
	Importante notare:
	\begin{itemize}
		\item A livello concettuale, l’unico attributo dell’entità Esami è Voto.\\
		Gli attributi di Esami che servono come riferimenti, cioè Studente (verso la tabella Studenti) e Corso (verso la tabella Corsi), sono attributi di Esami nel modello logico solo per tradurre l’informazione che nel diagramma E/R è fornita dal legame grafico e specifica che Esami è un’associazione che collega le entità Corsi e Studenti.
		\item Come attributi, Studente e Corso sono definiti su un dominio semplice (cioè ogni attributo ha un solo valore) e quindi fanno riferimento a una sola istanza di Studenti o Corsi.
	\end{itemize}
	\subsubsection{Esempio}
	\begin{tabular}{ |c|c|c|c| }
		\hline
		\multicolumn{4}{|c|}{Studenti}\\
		\hline
		Matricola & Cognome & Nome & Data di nascita\\
		\hline
		6554 & Rossi & Mario & 05/12/1978\\
		8765 & Neri & Paolo & 03/11/1976\\
		9283 & Verdi & Luisa & 12/11/1979\\
		3456 & Rossi & Maria & 01/02/1978\\
		\hline
	\end{tabular}
	\begin{tabular}{ |c|c|c| }
		\hline
		\multicolumn{3}{|c|}{Esami}\\
		\hline
		Studente & Voto & Corso\\
		\hline
		3456 & 30 & 04\\
		3456 & 24 & 01\\
		9283 & 28 & 01\\
		6554 & 26 & 02\\
		\hline
	\end{tabular}
	\begin{tabular}{ |c|c|c| }
		\hline
		\multicolumn{3}{|c|}{Corsi}\\
		\hline
		Codice & Titolo & Docente\\
		\hline
		01 & Analisi & Mario\\
		02 & Chimica & Bruni\\
		04 & Chimica & Verdi\\
		\hline
	\end{tabular}
	\begin{figure}[h]
		\centering
		\adjustbox{cfbox=white 5pt 0cm}{\includegraphics[scale=0.3]{example-image-golden}}
		\caption{Esplicitazione relazioni fra le tabelle}
		\label{fig:im-7}
	\end{figure}
	Dall’esempio abbiamo visto che:
	\begin{itemize}
		\item Il modello relazionale è basato su valori.
		\item I riferimenti fra dati in relazioni diverse avvengono attraverso la corrispondenza dei valori con i quali (insiemi di) attributi corrispondenti vengono istanziati nelle tuple che sono logicamente correlate.
	\end{itemize}
	NB Attributi corrispondenti può implicare uguaglianza del nome oppure corrispondenza (logica) fra valori ammissibili, attraverso l’imposizione di vincoli. In genere i due (insiemi di) attributi ammettono gli stessi valori oppure, per uno dei due, sono ammissibili i valori di un sottoinsieme dei valori ammissibili per l’altro.\\
	Gli altri modelli (gerarchico, reticolare) utilizzano puntatori per realizzare la corrispondenza.
	\subsubsection{Vantaggi dell’approccio basato su valori}
	\begin{itemize}
		\item Si inseriscono nella base di dati solo valori significativi per l’applicazione (i puntatori sono dati aggiuntivi relativi alla sola implementazione fisica).
		\item Il trasferimento dei dati da un ambiente ad un altro è più semplice (i puntatori hanno validità solo locale)
		\item La rappresentazione logica dei dati non fa riferimento a quella fisica e quindi si ottiene l’indipendenza (fisica) dei dati
	\end{itemize}
	\subsection{Informazione Sbagliata}
	Le tuple che compongono la base di dati devono essere omogenee.\\
	Quindi, in ogni tupla, ad ogni attributo deve essere associato un valore.\\
	Non sempre questo è possibile.\\ \\
	Es. Persone($Cognome$, $Nom$e, $Indirizzo$, $Telefono$)\\ \\
	Potrebbe esistere una persona che non possiede telefono, o di cui non conosciamo l’indirizzo.
	\begin{center}
		\begin{tabular}{ |c|c|c|c| }
			\hline
			Cognome & Nome & Indirizzo & Telefono\\
			\hline
			Rossi & Francesco & v. Rossa 22 & 0521 335643\\
			Verdi & Giuseppe & v. Verde 11 & \\
			Bruni & Bruno & &336 7564213\\
			\hline
		\end{tabular}
	\end{center}
	Questa relazione non è ammissibile! In ogni tupla, ogni attributo deve essere istanziato.\\
	Non conviene (anche se è un espediente di uso comune) usare valori del dominio normalmente non utilizzati (0, stringa nulla, “99”, ...), come spesso accade nella programmazione:
	\begin{itemize}
		\item Potrebbero non esistere valori non utilizzati
		\item Valori non utilizzati potrebbero, a un certo punto, diventare significativi ed essere utilizzati
		\item In fase di utilizzo (nei programmi) sarebbe necessario ogni volta tenere conto del “significato” (non standardizzato) di questi valori
	\end{itemize}
	Nel modello relazionale è definito un valore convenzionale, detto valore nullo, che indica la non disponibilità dell’informazione.\\
	Il valore nullo può rappresentare correttamente tre situazioni logicamente diverse in cui l’informazione è:
	\begin{itemize}
		\item Sconosciuta (so che il valore esiste ma non lo conosco)
		\item Inesistente (so che il valore non esiste)
		\item Indeterminata (non so se il valore esiste e, in ogni caso, non lo conosco)
	\end{itemize}
	\subsection{Vincoli di Integrità}
	Non tutte le combinazioni possibili di valori dei domini su cui è definita una relazione sono accettabili.
	\begin{itemize}
		\item Alcuni attributi possono assumere valori solo in un certo intervallo
		\item Alcuni attributi devono essere diversi in ogni tupla della stessa relazione
	\end{itemize}
	Es. valori dell’attributo Matricola nella relazione\\
	Studenti($Matricola$,$Cognome$,$Nome$,$DataNascita$)\\ \\
	Alcuni valori possono essere incompatibili con altri all’interno della stessa relazione\\ \\
	Es. data la relazione \\
	Esami($Matricola$,$Voto$,$Lode$,$CodCorso$)
	\begin{enumerate}
		\item Una stessa coppia Matricola, Corso può apparire una sola volta
		\item Il valore Vero per l’attributo Lode è corretto solo se Voto=30
	\end{enumerate}
	Alcuni valori possono essere incompatibili con i valori di un’altra relazione\\ \\
	Es. date le relazioni:\\
	Studenti ($Matricola$, $Cognome$, $Nome$, $DataNascita$)\\
	Esami ($Studente$, $Voto$, $Lode$, $CodCorso$)\\
	Corsi ($CodCorso$, $Titolo$, $Docente$) \\ \\
	\begin{enumerate}
		\item Ogni valore di CodCorso in Esami, in questo caso, DEVE essere un valore esistente di CodCorso nella relazione Corsi;
		\item Analogamente, ogni valore dell’attributo Studente nella tabella Esami DEVE essere un valore esistente dell’attributo Matricola nella relazione Studenti.
	\end{enumerate}
	Sono condizioni, espresse come predicati logici, inserite nella base di dati per garantirne la consistenza.\\
	Ogni istanza della base di dati deve soddisfare i vincoli di integrità, cioè il predicato corrispondente al vincolo deve assumere valore vero \textbf{in ogni istante}.\\
	Una istanza che soddisfi tutti i vincoli è detta corretta (o lecita, o ammissibile)\\ \\
	Un vincolo è detto:
	\begin{itemize}
		\item \textbf{Intra-relazionale} se coinvolge attributi della stessa relazione
	\end{itemize}
	Es.
	\begin{itemize}
		\item \textbf{Vincoli di tupla} possono essere valutati su ciascuna tupla indipendentemente dalle altre
		\item \textbf{Vincoli di dominio} definiti su singoli valori
		\item \textbf{Interrelazionale} se coinvolge più relazioni
	\end{itemize}
	\subsubsection{Vincoli di Tupla}
	Possono essere definiti attraverso operatori booleani\\
	Es. Data la relazione:\\
	Esami($Matricola$,$Voto$,$Lode$,$CodCorso$)
	\begin{lstlisting}[language=SQL]
		(Voto >= 18)AND(Voto <= 30) 
		(NOT (lode=Vero)) OR (Voto=30)
	\end{lstlisting}
	Oppure, data la relazione:\\
	Pagamenti($Data$,$Importo$,$Ritenute$,$Netto$)
	\begin{lstlisting}[language=SQL]
		Netto = Importo - Ritenute
	\end{lstlisting}
	\subsubsection{Identificazione delle tuple}
	\begin{center}
		\begin{tabular}{ |c|c|c|c|c| }
			\hline
			\textbf{Matricola} & \textbf{Cognome} & \textbf{Nome} & \textbf{Corso} & \textbf{Nascita}\\
			\hline
			27655 & Rossi & Mario & Ing Inf & 5/12/78\\
			78763 & Rossi & Mario & Ing Inf & 3/11/76\\
			65432 & Neri & Piero & Ing Mecc & 10/7/79\\
			87654 & Neri & Mario & Ing Inf & 3/11/76\\
			67653 & Rossi & Piero & Ing Mecc & 5/12/78\\
			\hline
		\end{tabular}
	\end{center}
	Osservazioni:
	\begin{itemize}
		\item Non ci sono due tuple con lo stesso valore sull’attributo Matricola
		\item Non ci sono due tuple uguali su tutti e tre gli attributi Cognome, Nome e Nascita
	\end{itemize}
	\subsection{Chiavi}
	Una chiave è un insieme minimale di attributi utilizzato per identificare univocamente le tuple di una relazione.\\
	\textbf{Formalmente}: Un insieme di attributi $K$ è superchiave per una relazione $r$ se $r$ non contiene due tuple $t_1$ e $t_2$ tali che  $t_1[K] = t_2[K]$\\
	Un insieme di attributi $K$ è chiave per $r$ se è superchiave minimale, cioè se non esiste un’altra superchiave $K'$ che sia sottoinsieme di $K$ o, comunque, di grado inferiore.\\ \\
	Una chiave è tale se soddisfa la definizione per tutte le possibili tuple appartenenti alla relazione, e non solo per quelle che effettivamente appaiono come istanze della relazione stessa.\\
	$\Rightarrow$La chiave è legata allo schema della relazione e non ai valori effettivamente assunti dalle istanze dello schema.\\
	Ogni relazione, per definizione, possiede una chiave.\\
	Infatti, poiché una relazione non ammette due tuple uguali, l’intero insieme di attributi $X$ su cui è definita è sicuramente superchiave per la relazione.\\ \\
	Può possederne anche più di una.
	\subsubsection{Esempio}
	\begin{center}
		\begin{tabular}{ |c|c|c|c|c| }
			\hline
			\textbf{Matricola} & \textbf{Cognome} & \textbf{Nome} & \textbf{Corso} & \textbf{Nascita}\\
			\hline
			27655 & Rossi & Mario & Ing Inf & 5/12/78\\
			78763 & Rossi & Mario & Ing Inf & 3/11/76\\
			65432 & Neri & Piero & Ing Mecc & 10/7/79\\
			87654 & Neri & Mario & Ing Inf & 3/11/76\\
			67653 & Rossi & Piero & Ing Mecc & 5/12/78\\
			\hline
		\end{tabular}
	\end{center}
	Matricola è una chiave:
	\begin{itemize}
		\item è superchiave
		\item contiene un solo attributo e quindi è minimale
	\end{itemize}
	\subsubsection{Chiave primaria}
	La presenza di valori nulli in una chiave può vanificare la proprietà di unicità delle tuple che identifica.\\
	Si impone quindi che almeno una chiave non contenga valori nulli.\\
	Tale chiave è detta chiave primaria.\\
	Di solito la chiave primaria compare sottolineata nello schema di una relazione.\\
	Es.\\
	Studenti($Matricola$, $Cognome$, $Nome$, $Nascita$, $Corso$)
	\subsubsection{Vincoli si Integrità Referenziale}
	In alcuni casi (corrispondenze fra relazioni) è necessario che i valori degli attributi di una relazione $R_1$ si trovino anche in attributi corrispondenti di un’altra relazione $R_2$.\\
	Un \textbf{vincolo di integrità referenziale} (o \textbf{foreign key}, o \textbf{chiave esterna}) fra un insieme di attributi $X$ di $R_1$ e un’altra relazione $R_2$ è soddisfatto se i valori su $X$ di ciascuna tupla di $R_1$ compaiono come valori della chiave (primaria) di $R_2$.
	\paragraph{Esempio}
	Corsi ($Corso$, $Ncrediti$)\\
	Manifesto ($CdL$, $Insegnamento$, $Anno$)\\
	Esiste un vincolo di integrità referenziale fra Insegnamento, attributo di Manifesto, e Corso, chiave primaria di Corsi
	\begin{center}
		\begin{tabular}{ |c|c| }
			\hline
			\textbf{Corso} & \textbf{Ncrediti}\\
			\hline
			Basi di Dati e Web & 9\\
			Sistemi Operativi & 6\\
			Ingegneria Del Software & 9\\
			\hline
		\end{tabular}
		\hfill
		\begin{tabular}{ |c|c|c| }
			\hline
			\textbf{CdL} & \textbf{Insegnamento} & \textbf{Anno}\\
			\hline
			IET-I & Basi di Dati e Web & 3\\
			IET & Sistemi Operativi & 2\\
			IET-I & Ingegneria del Software & 3\\
			\hline
		\end{tabular}
	\end{center}
\end{document}